% ============================================================================
% Fichier  : partie3/P03_C12_droit_camerounais_africain.tex
% Titre    : Droit Camerounais et Africain
% Auteur   : MINKA MI NGUIDJOÏ Thierry Emmanuel
% Version  : 1.0 -- Juin 2026
% Statut   : RÉDIGÉ
% Sources  : M3_Cadre_Normatif_National_PNC.docx (intégral -- textes articles,
%            tableaux sanctions, règles articulation, simulation MOABI)
%            14_droit_camerounais_africain.tex (squelette 98 lignes)
%            Prologue (affaire MBONGO), Ch. custody (art. 52-59)
% Positionnement : Ce chapitre est le plus important du cours pour la pratique
%                  quotidienne des OPJ camerounais : qualification correcte,
%                  rédaction de PV, transmissibilité au parquet.
%                  Le M3 PNC est la source principale -- fidèlement transposé.
% ============================================================================

\chapter{Droit Camerounais et Africain}
\label{chap:droit-cam}

\epigraph{%
    « Africa must develop its own digital legal framework that respects
    both international standards and local cultural realities. »%
}{--- Nnenna Ifeanyi-Ajufo}

\niveauChapitre{Fondamental}{%
    Notions de droit général recommandées mais non indispensables~:
    ce chapitre part de zéro sur le droit camerounais du numérique.%
}

\begin{boxobjectifs}
\begin{itemize}
    \item Identifier les trois textes du cadre normatif national et
          leur articulation~: \loicam{2010/012}, \loicam{2024/017},
          Code pénal.
    \item Qualifier les infractions numériques les plus fréquentes
          au Cameroun (fraude, accès illicite, ransomware, CSAM) selon
          l'article de loi exact.
    \item Appliquer la \loicam{2024/017} aux violations de données
          personnelles et identifier les sanctions applicables.
    \item Maîtriser les règles d'articulation des trois textes~:
          quand cumuler, quand choisir.
    \item Identifier les six erreurs de qualification les plus fréquentes
          et leurs conséquences sur la validité du dossier judiciaire.
    \item Rédiger un rapport de qualification d'infraction numérique
          transmissible au parquet.
\end{itemize}
\end{boxobjectifs}

\bigskip

% ============================================================================
\section*{Ce chapitre dans la pratique quotidienne}
% ============================================================================

Un OPJ qui ne maîtrise pas les qualifications nationales sera incapable
de rédiger un procès-verbal valable, de choisir le bon article d'inculpation
et de préparer un dossier transmissible au parquet. Les droits internationaux
des chapitres précédents sont inutilisables sans cette fondation nationale.

\textbf{Méthode de ce chapitre~:} partir systématiquement de cas concrets
camerounais~--- fraudes MoMo, escroqueries Facebook, SIM~swapping, ransomware~---
et remonter vers la qualification légale. Jamais l'inverse.

\separateur

% ============================================================================
\section{Architecture du Cadre Normatif National~: Trois Textes}
\label{sec:architecture-nationale}
% ============================================================================

\subsection{Vue d'ensemble}

Trois textes constituent le droit camerounais du numérique~:

\begin{tableaucours}{p{3.5cm}p{4.0cm}p{4.5cm}}{%
    \tableauHeader{Texte} &
    \tableauHeader{Objet} &
    \tableauHeader{Rôle pour l'enquêteur}
}
\textbf{\loicam{2010/012}} du 21~déc.~2010 &
    Cybersécurité et cybercriminalité &
    \textbf{Texte principal de qualification pénale.}
    97~articles, 5~titres, 73~définitions. \\
\textbf{\loicam{2024/017}} du 23~déc.~2024 &
    Protection des données à caractère personnel &
    Complète la Loi~2010 pour les violations de données.
    Sanctions administratives ET pénales. \\
\textbf{Code pénal camerounais} &
    Droit pénal général &
    S'applique en cumul pour l'escroquerie, le faux,
    la diffamation, les menaces. \\
\end{tableaucours}
\vspace{-0.8em}
\begin{center}{\small\itshape Tableau~12.1 --- Les trois textes du cadre national}\end{center}

\begin{boxalert}
\textbf{Quatre points clés avant de qualifier~:}
\begin{enumerate}
    \item L'article~4 de la Loi~2010/012 définit \textbf{73~termes}. Ces
          définitions ont valeur juridique contraignante. Une qualification
          fondée sur une définition incorrecte fragilise le dossier.
    \item La Loi~2010 n'abroge pas le Code pénal. \textbf{Les deux se cumulent.}
          Un même fait peut être à la fois escroquerie (CP) et fraude
          informatique (art.~72 Loi~2010).
    \item \textbf{Art.~89 Loi~2010~: le sursis est impossible.} Tout condamné
          sous la Loi~2010 purge sa peine. Mentionner ce point dans
          chaque rapport transmis au parquet.
    \item Le champ d'application est large~: réseaux + systèmes d'information
          + équipements terminaux. \textbf{Un téléphone mobile est un
          équipement terminal. Un groupe WhatsApp est un réseau de
          communications électroniques.}
\end{enumerate}
\end{boxalert}

\subsection{Architecture de la Loi~2010/012}

\begin{tableaucours}{p{3.5cm}p{8.0cm}}{%
    \tableauHeader{Titre} &
    \tableauHeader{Contenu et articles clés}
}
\textbf{Titre~I} --- Dispositions générales &
    Définitions (art.~4 --- 73~termes).
    Champ d'application. Exclusion défense/sécurité nationale (art.~2). \\
\textbf{Titre~II} --- Cybersécurité &
    Politique de sécurité (art.~6). Régulation ANTIC (art.~7--9).
    Certification et signature électronique (art.~10--23).
    Protection des réseaux et vie privée (art.~24--51). \\
\textbf{Titre~III} --- Cybercriminalité &
    \textbf{Dispositions processuelles (art.~52--59)~: perquisitions,
    saisies, réquisitions.}
    \textbf{Infractions et sanctions (art.~60--89)~: toutes les
    qualifications pénales.} \\
\textbf{Titre~IV} --- Coopération judiciaire &
    Entraide judiciaire internationale (art.~91--94). \\
\textbf{Titre~V} --- Dispositions finales &
    Textes d'application (art.~95--97). \\
\end{tableaucours}
\vspace{-0.8em}
\begin{center}{\small\itshape Tableau~12.2 --- Architecture de la Loi~2010/012}\end{center}

% ============================================================================
\section{Les Infractions de la Loi~2010/012~: Titre~III}
\label{sec:infractions-loi2010}
% ============================================================================

\subsection{Articles~65 à~69~: accès, interception, atteinte à l'intégrité}

\begin{boxjuris}
\textbf{Article~65 --- Accès illicite et interception illégale~:}
\begin{quote}
(1)~Est puni d'un emprisonnement de cinq~(05) à dix~(10)~ans et d'une
amende de 5~000~000 à 10~000~000~F~CFA \textit{[\ldots]} celui qui effectue,
sans droit ni autorisation, l'interception par des moyens techniques, de
données lors des transmissions \textit{[\ldots]} à destination, en provenance
ou à l'intérieur ou non d'un réseau de communications électroniques.

(2)~Est puni des mêmes peines tout accès non autorisé à l'ensemble ou à
une partie d'un réseau de communications électroniques ou d'un système
d'information ou d'un équipement terminal.

(3)~Les peines prévues à l'alinéa~1 ci-dessus sont doublées, en cas
d'accès illicite portant atteinte à l'intégrité, la confidentialité,
la disponibilité du réseau ou du système d'information.

(4)~Est puni des mêmes peines celui qui, sans droit, permet l'accès dans
un réseau ou dans un système d'information par défi intellectuel.
\end{quote}
\textbf{L'article fondamental.} Toute intrusion informatique non autorisée
tombe sous cet article --- qu'elle ait causé un dommage ou non. L'argument
de la ``simple curiosité'' ne constitue pas une défense~: l'al.~4 le
criminalise explicitement.\\[0.3ex]
\textbf{Circonstance aggravante (al.~3)~:} peines doublées si l'accès porte
atteinte à l'intégrité, la confidentialité ou la disponibilité.
En pratique~: toute intrusion qui modifie, copie ou rend inaccessibles
des données tombe dans l'aggravation.\\[0.3ex]
\textbf{Cas type PNC~:} accès au compte bancaire en ligne d'un tiers en
utilisant ses identifiants volés~$\to$ art.~65~al.~2 + art.~65~al.~3
(données copiées).
\end{boxjuris}

\begin{boxjuris}
\textbf{Articles~66--67 --- Atteinte à l'intégrité des réseaux et systèmes~:}
Art.~66~al.~1~: puni de 2 à 5~ans + 1~000~000 à 2~000~000~F~CFA, celui
qui entraîne la perturbation ou l'interruption du fonctionnement d'un
réseau en introduisant, transmettant, endommageant, effaçant, détériorant,
modifiant, supprimant ou rendant inaccessibles les données.\\[0.3ex]
\textbf{Distinction art.~66 vs art.~67~:} art.~66~= perturbation d'un réseau
ou équipement terminal~; art.~67~= atteinte à l'intégrité d'un
système d'information. Les deux peuvent être cumulés.\\[0.3ex]
\textbf{Couverture pratique~:} attaques DDoS (saturation~= perturbation)~;
ransomware (chiffrement~= rendre inaccessibles)~; defacement de site web
(modification~= art.~67)~; sabotage de base de données.
\end{boxjuris}

\begin{boxjuris}
\textbf{Articles~68--69 --- Accès aggravé et espionnage informatique~:}
Art.~68~: puni de 5 à 10~ans + 10~M à 50~M~F~CFA celui qui accède ou
se maintient frauduleusement dans un système en causant une perturbation.
\textbf{Peines doublées si suppression ou modification de données
(art.~68~al.~2).}\\[0.3ex]
Art.~69~: accès en violation des mesures de sécurité pour obtenir des
données d'un système interconnecté. Amende plafond~: \textbf{100~millions
FCFA.}\\[0.3ex]
\textbf{Cas type art.~69~:} intrusion dans le système de paie d'une
administration camerounaise depuis Internet pour extraire des données
salariales.
\end{boxjuris}

\subsection{Articles~70 à~72 et~86~: fraude, manipulation, abus de dispositifs}

\begin{boxjuris}
\textbf{Article~72 --- Fraude informatique~:}
Est puni des peines de l'art.~66 celui qui, de quelque manière que ce
soit, sans droit, introduit, altère, efface ou supprime, \textbf{afin
d'obtenir un bénéfice économique}, les données électroniques, de manière
à causer un préjudice patrimonial à autrui.\\[0.3ex]
\textbf{Trois éléments constitutifs~:} (1)~acte sur des données
(introduction, altération, effacement, suppression), (2)~sans droit,
(3)~pour obtenir un bénéfice économique causant un préjudice patrimonial.\\[0.3ex]
\textbf{L'article le plus utilisé pour les escroqueries en ligne.}
BEC, phishing bancaire, fraude MoMo, modification de relevés bancaires~---
tous tombent sous art.~72.\\[0.3ex]
\textbf{Cumul recommandé~:} art.~72 Loi~2010 + art.~318 Code pénal
(escroquerie) quand la victime a été manipulée ET que des données
ont été altérées.
\end{boxjuris}

\begin{boxjuris}
\textbf{Article~86 --- Abus de dispositifs informatiques~:}
Est puni d'un emprisonnement de 2 à 5~ans \textit{[\ldots]} celui qui
importe, détient, offre, cède, vend ou met à disposition un programme
informatique, un mot de passe, un code d'accès ou toutes données
informatiques similaires conçus et/ou spécialement adaptés pour permettre
d'accéder à un système d'information.\\[0.3ex]
\textbf{Criminalise la préparation~:} vendre un kit de phishing, partager
des identifiants volés, mettre à disposition un logiciel d'intrusion~---
même sans attaque effective.\\[0.3ex]
\textbf{Peines accessoires (art.~87)~:} confiscation des outils~;
interdiction d'exercer une activité professionnelle (5~ans)~;
fermeture d'établissement~; exclusion des marchés publics.
\end{boxjuris}

\subsection{Article~88~: refus de remettre la clé de déchiffrement}

\begin{boxjuris}
\textbf{Article~88 --- Refus de déchiffrement~:}
(1)~Est puni d'un emprisonnement de \textbf{un~(01) à cinq~(05)~ans} et
d'une amende de \textbf{100~000 à 1~000~000~F~CFA} celui qui refuse de
remettre aux autorités judiciaires compétentes la convention secrète de
déchiffrement d'un moyen cryptologique susceptible d'avoir été utilisé
pour préparer ou commettre une infraction.

(2)~Les peines prévues à l'alinéa~1 sont portées de 3 à 5~ans et de
1~000~000 à 5~000~000~F~CFA lorsque la remise de ladite convention aurait
permis d'éviter un crime ou délit ou d'en limiter les effets.
\end{boxjuris}

\begin{boxPNC}
\textbf{L'article 88~: réflexe obligatoire face à un téléphone chiffré}

Un suspect refuse de déverrouiller son téléphone. L'OPJ ne doit pas bloquer~:
\begin{enumerate}
    \item Obtenir une réquisition formelle du procureur ou du juge
          d'instruction (la réquisition est la condition préalable).
    \item Si le suspect refuse après réquisition~: art.~88~al.~1
          constitue une infraction autonome \emph{indépendante} de
          l'infraction principale.
    \item Si des éléments montrent que le déchiffrement aurait permis
          d'éviter un crime en cours~: art.~88~al.~2 (aggravé~: 3--5~ans).
    \item Documenter précisément~: date et heure de la réquisition,
          formulation exacte de la demande, refus exprimé clairement.
\end{enumerate}
\textbf{Erreur fréquente~:} beaucoup d'enquêteurs ignorent l'existence
de cet article. Un dossier CSAM ou ransomware peut être débloqué par
la seule citation de l'art.~88.
\end{boxPNC}

\subsection{Article~76~: CSAM et contenus sexuels impliquant des mineurs}

\begin{boxjuris}
\textbf{Article~76 --- Pornographie enfantine~:}
Est puni d'un emprisonnement de \textbf{cinq~(05) à dix~(10)~ans} et
d'une amende de \textbf{5~000~000 à 10~000~000~F~CFA} celui qui fixe,
enregistre, transmet ou diffuse, via un système d'information ou un
réseau de communications électroniques, des images pornographiques
impliquant des mineurs.\\[0.3ex]
\textbf{Ressource opérationnelle clé~:} NCMEC CyberTipline
(\texttt{cybertipline.org}). Les fournisseurs américains ont l'obligation
légale de signaler tout CSAM détecté~; les signalements sont transmis
via INTERPOL.\\[0.3ex]
\textbf{Cumul obligatoire~:} art.~76 Loi~2010 + art.~295 Code pénal
(attentat à la pudeur aggravé) + art.~80--81 Loi~2010 si proposition
sexuelle en ligne à un mineur.
\end{boxjuris}

% ============================================================================
\section{La Loi~2024/017~: Protection des Données Personnelles}
\label{sec:loi2024}
% ============================================================================

\subsection{Architecture et innovations par rapport à la Loi~2010}

La \loicam{2024/017} du 23~décembre~2024 comble un vide~: la Loi~2010
réprimait les violations de données, mais ne fixait pas les obligations
préventives (consentement, sécurité obligatoire, registre de traitement,
notification de violation). La Loi~2024 établit un système complet
analogue au RGPD européen.

\begin{tableaucours}{p{3.5cm}p{4.0cm}p{4.0cm}}{%
    \tableauHeader{Titre / Chapitre} &
    \tableauHeader{Objet} &
    \tableauHeader{Articles clés}
}
Titre~I --- Général & Champ d'application, définitions fondamentales &
    Art.~1--5 \\
Titre~II --- Principes & Légitimité, consentement, finalité, minimisation,
    sécurité &
    Art.~6--36 \\
Titre~III --- Droits & Accès, rectification, effacement, portabilité,
    opposition &
    Art.~37--47 \\
Titre~IV --- Interdictions & Données sensibles, profilage, traitements
    illicites &
    Art.~48--52 \\
Titre~V --- Autorité & Création et pouvoirs de l'Autorité de protection &
    Art.~53 \\
Titre~VI --- Sanctions & Administratives + pénales &
    Art.~54--71 \\
\end{tableaucours}
\vspace{-0.8em}
\begin{center}{\small\itshape Tableau~12.3 --- Architecture de la Loi~2024/017}\end{center}

\subsection{Définitions fondamentales~: art.~5}

\begin{itemize}
    \item \textbf{Donnée personnelle~:} toute information permettant
          d'identifier directement ou indirectement une personne physique.
          \textbf{L'adresse IP est une donnée personnelle.} Une intrusion
          qui consigne les adresses IP des victimes active la Loi~2024.
    \item \textbf{Donnée sensible~:} origine tribale ou raciale, opinions
          politiques, convictions religieuses, appartenance syndicale,
          données de santé, vie sexuelle, données génétiques, biométriques.
          \textbf{Protection renforcée~: peines doublées (art.~74~al.~5
          Loi~2010, art.~66 Loi~2024).}
    \item \textbf{Responsable du traitement~:} toute personne physique
          ou morale qui détermine les finalités et les moyens du traitement.
    \item \textbf{Profilage~:} traitement automatisé visant à évaluer des
          aspects personnels relatifs à une personne (santé, préférences,
          localisation, situation économique).
\end{itemize}

\subsection{Obligations clés du responsable du traitement}

\begin{boxjuris}
\textbf{Article~22 --- Notification de violation de données (``breach'')}\\
Le responsable du traitement notifie \textbf{sans délai} à l'Autorité de
protection toute violation de données à caractère personnel susceptible
d'engendrer un risque pour les droits et libertés des personnes concernées.

\textbf{Sanction du défaut de notification~:} art.~56~: amende
administrative de 1 à 10~millions FCFA.\\[0.3ex]
\textbf{Pour l'enquêteur~:} quand une entreprise victime d'un ransomware
ou d'une fuite de données n'a pas notifié l'Autorité~--- noter ce manquement
dans le rapport d'enquête. C'est une infraction administrative indépendante,
sanctionnable en plus des infractions commises par l'attaquant.
\end{boxjuris}

\begin{boxjuris}
\textbf{Article~27 --- Sécurité obligatoire des données}\\
Le responsable du traitement met en \oe uvre \textit{[\ldots]} les mesures
techniques et organisationnelles appropriées~: \textbf{traçabilité des accès}
(journaux de connexion, logs d'audit)~; \textbf{contrôle d'accès}
(principe du moindre privilège)~; chiffrement des données sensibles~;
sauvegardes régulières et plans de continuité.

\textbf{Points utiles pour l'enquête~:} (1)~le responsable a l'obligation
de constituer et conserver des logs d'audit~--- ils existent, ils peuvent
être requis~; (2)~un employé qui accède à des données hors de son périmètre
commet une infraction ET son employeur est en faute s'il n'avait pas mis
en place les contrôles.
\end{boxjuris}

\begin{boxjuris}
\textbf{Article~32 --- Transfert international de données~:}
Le transfert de données à caractère personnel vers un État étranger
est subordonné à une \textbf{autorisation préalable} de l'Autorité
de protection, qui doit s'assurer~: que le pays de destination offre
un niveau de protection suffisant~; de l'entrée en vigueur d'un instrument
juridique avec ce pays~; de la souscription de clauses contractuelles types.\\[0.3ex]
\textbf{Sanction~: art.~69~: 3 à 10~ans + 2 à 20~millions FCFA.}
L'une des infractions les plus graves de la Loi~2024.
\end{boxjuris}

\subsection{Les sanctions pénales de la Loi~2024/017}

\begin{boxjuris}
\textbf{Article~63 --- Collecte frauduleuse de données personnelles~:}
(1)~Puni de \textbf{2 à 5~ans} + \textbf{200~000 à 5~000~000~F~CFA}
celui qui collecte des données à caractère personnel ou y accède par
un moyen frauduleux, déloyal ou illicite.

(2)~\textbf{Peines doublées} lorsque la collecte frauduleuse s'accompagne
d'un verrouillage, d'un cryptage ou de toute autre technique portant atteinte
à la disponibilité et à l'intégrité des données.

\textbf{Al.~2 couvre directement les ransomwares~:} l'attaque ransomware
= collecte frauduleuse + chiffrement~= peines doublées.\\[0.3ex]
\textbf{Cumul possible avec art.~65 Loi~2010} quand la collecte frauduleuse
implique une intrusion dans un système d'information.
\end{boxjuris}

% ============================================================================
\section{Tableau Comparatif des Sanctions~: les Trois Textes}
\label{sec:sanctions-comparees}
% ============================================================================

\begin{tableaucours}{p{3.5cm}p{2.8cm}p{2.8cm}p{2.5cm}}{%
    \tableauHeader{Infraction} &
    \tableauHeader{Loi~2010/012} &
    \tableauHeader{Loi~2024/017} &
    \tableauHeader{Code pénal (cumul)}
}
Accès illicite & Art.~65~: 5--10~ans + 5--10M & Art.~63 si données perso. &
    Tentative possible (art.~97) \\
Interception illégale & Art.~65~al.~1~: 5--10~ans & Art.~63 si données perso. &
    Art.~305 si contenu utilisé \\
Atteinte intégrité syst. & Art.~66--67~: 2--5~ans + 1--2M &
    Art.~63~al.~2 si chiffrement &
    Art.~318 si préjudice \\
Fraude informatique & Art.~72~: 2--5~ans + 1--2M &
    Art.~63 si données perso. &
    Art.~318 escroquerie \\
Collecte frauduleuse & Art.~74~al.~4~: 6~mois--2~ans + 1--5M &
    Art.~63~: 2--5~ans (plus lourd) & Art.~318 si revente \\
Données sensibles &
    Art.~74~al.~5~: peines~$\times$~2 &
    Art.~66~: 6~mois--2~ans + 500K--5M & Art.~305 \\
Transfert illégal étranger &
    Art.~74~al.~7~: 6~mois--2~ans + 5--50M &
    Art.~69~: \textbf{3--10~ans + 2--20M} & Art.~99 si données étatiques \\
Profilage sans consentement & Non prévu expressément &
    Art.~65~: 3--10~ans + 1--20M & Non prévu \\
CSAM & Art.~76/80~: 5--10~ans + 5--10M &
    Art.~52 atteinte à la dignité & Art.~295 aggravé \\
Refus déchiffrement & Art.~88~: \textbf{1--5~ans + 100K--1M} &
    Non prévu & Non prévu \\
Personnes morales & Art.~64~: 5--50M +~peine access. &
    Art.~71~: \textbf{50M à 1~milliard} & Dissolution possible \\
\end{tableaucours}
\vspace{-0.8em}
\begin{center}{\small\itshape Tableau~12.4 --- Comparatif des sanctions~:
valeurs en millions FCFA sauf mention}\end{center}

\begin{boiteRemarque}{Trois chiffres à mémoriser}
\begin{itemize}
    \item \textbf{100~millions FCFA} (Loi~2010~: espionnage, art.~69)
          --- l'amende la plus lourde de la Loi~2010.
    \item \textbf{1~milliard FCFA} (Loi~2024, art.~71~: personnes morales)
          --- l'amende la plus lourde du dispositif national.
    \item \textbf{Sursis impossible} sous Loi~2010 (art.~89) --- unique
          dans le droit camerounais.
\end{itemize}
\end{boiteRemarque}

% ============================================================================
\section{Règles d'Articulation des Trois Textes}
\label{sec:articulation}
% ============================================================================

\begin{tableaucours}{p{3.2cm}p{8.5cm}}{%
    \tableauHeader{Règle} &
    \tableauHeader{Contenu et exemples}
}
\textbf{Règle~1 --- Spécialité} &
    La Loi~2010 est la \textit{lex specialis} du droit pénal cyber.
    Elle prime sur le Code pénal pour les infractions qu'elle définit
    expressément. Mais le Code pénal s'applique en complément pour les
    éléments non couverts ou les qualifications aggravées. \\
\textbf{Règle~2 --- Cumul} &
    Un même fait peut être qualifié sous plusieurs textes si les éléments
    constitutifs sont indépendants. Loi~2010 + Code pénal~: toujours possible.
    Loi~2010 + Loi~2024~: possible si la violation porte à la fois sur
    l'intégrité du système ET sur les données personnelles. \\
\textbf{Règle~3 --- Plus solide} &
    Quand deux qualifications sont possibles, choisir celle dont les
    éléments constitutifs sont les plus faciles à prouver, pas
    nécessairement la peine la plus élevée. \\
\textbf{Règle~4 --- Données perso.} &
    Tout fait impliquant des données à caractère personnel (y compris
    l'adresse IP~: art.~5 Loi~2024) active potentiellement la Loi~2024
    EN PLUS des autres textes. \\
\textbf{Règle~5 --- Personnes morales} &
    La Loi~2024 prévoit des amendes pour personnes morales bien
    supérieures à la Loi~2010 (jusqu'à 1~milliard FCFA vs 50~millions).
    Pour les infractions d'entreprises~: toujours vérifier si la Loi~2024
    est applicable. \\
\textbf{Règle~6 --- Sursis} &
    Art.~89 Loi~2010~: sursis impossible.
    Loi~2024~: pas de disposition similaire --- le sursis est en principe
    possible (Code pénal). Cette différence peut orienter la stratégie
    du parquet. \\
\end{tableaucours}
\vspace{-0.8em}
\begin{center}{\small\itshape Tableau~12.5 --- Règles d'articulation des trois textes}\end{center}

\subsection{Grille de qualification par type d'infraction courant}

\begin{tableaucours}{p{3.8cm}p{7.8cm}}{%
    \tableauHeader{Infraction type} &
    \tableauHeader{Articles à citer (toujours au moins deux textes)}
}
Fraude MoMo / BEC &
    Art.~72 Loi~2010 + Art.~318 CP (escroquerie)~;
    si données perso.~: +~Art.~63 Loi~2024 \\
Intrusion système &
    Art.~65~al.~2 Loi~2010~; aggravé~: art.~65~al.~3 si dommage~;
    + Art.~65 Loi~2024 si données perso. \\
Ransomware &
    Art.~66--67 Loi~2010 + Art.~63~al.~2 Loi~2024 (peines doublées)
    + Art.~22 Loi~2024 si pas de notification \\
Vol de base de données &
    Art.~63 Loi~2024 + Art.~74 Loi~2010 + Art.~318 CP (revente) \\
CSAM &
    Art.~76 + Art.~80--81 Loi~2010. Signaler NCMEC CyberTipline.
    + Art.~295 CP aggravé \\
Phishing / site frauduleux &
    Art.~72 + Art.~86 Loi~2010 + Art.~318 CP + Art.~155 CP (faux) \\
T�léphone chiffré, refus &
    Art.~88 Loi~2010. Réquisition judiciaire obligatoire préalablement. \\
Transfert données étranger &
    Art.~69 Loi~2024 (3--10~ans)~: la sanction la plus lourde du dispositif. \\
\end{tableaucours}
\vspace{-0.8em}
\begin{center}{\small\itshape Tableau~12.6 --- Grille de qualification par infraction type}\end{center}

% ============================================================================
\section{Six Erreurs de Qualification et Comment les Éviter}
\label{sec:erreurs-qualification}
% ============================================================================

\begin{tableaucours}{p{3.8cm}p{4.0cm}p{4.0cm}}{%
    \tableauHeader{Erreur} &
    \tableauHeader{Description} &
    \tableauHeader{Correction}
}
\textbf{1.~Tout qualifier sous art.~74 Loi~2010} &
    Art.~74 couvre la vie privée. Beaucoup d'enquêteurs l'utilisent
    pour toutes les infractions sur les données~: renvoi de dossier. &
    Pour les violations de données personnelles post-23~déc.~2024~:
    utiliser les articles pertinents Loi~2024. Art.~74 reste applicable
    pour les faits antérieurs. \\
\textbf{2.~Oublier art.~88} &
    Suspect refuse de déverrouiller son téléphone. Blocage total.
    L'enquêteur ignore qu'il existe une infraction autonome. &
    Réquisition du juge~$\to$ refus~= art.~88. Infraction indépendante
    de l'infraction principale. \\
\textbf{3.~Ne pas citer le Code pénal} &
    PV ne cite que Loi~2010. Le magistrat rejette la qualification pour
    un vice de forme. Pas de filet de sécurité. &
    Systématiquement citer en alternative~: art.~318~CP (escroquerie),
    art.~155--162 (faux), art.~305--308 (diffamation). \\
\textbf{4.~Confondre art.~66 et art.~72} &
    Art.~66~= perturbation sans bénéfice économique. Art.~72~= manipulation
    avec préjudice patrimonial. Qualifier une fraude sous art.~66 est faux. &
    Élément clé~: y a-t-il un bénéfice économique et un préjudice
    patrimonial~? Oui~$\to$ art.~72. Non (sabotage pur)~$\to$ art.~66. \\
\textbf{5.~Ignorer que l'IP est une donnée perso.} &
    OPJ note une IP sans penser à qualifier sous Loi~2024 la collecte
    de ces IP par un acteur non autorisé. &
    Art.~5 Loi~2024~: l'adresse IP est explicitement une donnée personnelle.
    Toute collecte sans base légale active la Loi~2024. \\
\textbf{6.~Ne pas noter le défaut de notification} &
    Entreprise victime d'un ransomware n'a pas notifié l'Autorité.
    L'OPJ ne le signale pas dans son rapport. &
    Art.~22 Loi~2024~= obligation de notifier ``sans délai''. Défaut~=
    infraction administrative indépendante (art.~56~: amende 1--10M). \\
\end{tableaucours}
\vspace{-0.8em}
\begin{center}{\small\itshape Tableau~12.7 --- Six erreurs de qualification et corrections}\end{center}

% ============================================================================
\section{Le Cadre Procédural~: De la Plainte au Jugement}
\label{sec:procedure}
% ============================================================================

\subsection{Autorités compétentes}

\begin{itemize}
    \item \textbf{ANTIC} (Agence Nationale des TIC)~: autorité de régulation
          technique~; point focal de la coopération internationale cyber~;
          compétences d'enquête administrative.
    \item \textbf{DGSN/DCPJ}~: Direction centrale de la Police Judiciaire~;
          BCN INTERPOL Cameroun~; unité de cybercriminalité.
    \item \textbf{Parquet spécialisé}~: Tribunal de Grande Instance de Yaoundé
          et Douala disposent de magistrats affectés aux affaires cyber.
    \item \textbf{Autorité de protection des données (Loi~2024, art.~53)}~:
          organisme public indépendant~; délivre les autorisations~;
          dispose d'un arsenal de sanctions administratives.
          \textbf{Statut~(juin~2026)~:} à créer par décret présidentiel~;
          surveiller le Journal Officiel.
    \item \textbf{Experts agréés}~: Décret~N°~69/DF/544~; BAC+5 informatique
          minimum~; 5~ans d'expérience~; agrément du Ministère de la Justice.
\end{itemize}

\subsection{Procédure type d'investigation numérique}

\begin{figure}[H]
\centering
\begin{tikzpicture}[
    etape/.style={draw=CouleurPrimaire, fill=CouleurDef, rounded corners=3pt,
        minimum width=7cm, minimum height=0.85cm,
        font=\small\bfseries\color{CouleurPrimaire}, align=center},
    fl/.style={-{Stealth[length=4pt]}, thick, color=CouleurSecondaire}
]
\node[etape] (e1) at (0, 0)   {1. Plainte / Signalement};
\node[etape] (e2) at (0,-1.1) {2. Enquête préliminaire (OPJ + ANTIC)};
\node[etape] (e3) at (0,-2.2) {3. Ouverture information judiciaire};
\node[etape] (e4) at (0,-3.3) {4. Commission rogatoire pour expertise};
\node[etape] (e5) at (0,-4.4) {5. Expertise technique (expert agréé)};
\node[etape] (e6) at (0,-5.5) {6. Rapport d'expertise};
\node[etape] (e7) at (0,-6.6) {7. Audience (présentation des preuves)};
\node[etape] (e8) at (0,-7.7) {8. Jugement};
\draw[fl] (e1)--(e2); \draw[fl] (e2)--(e3); \draw[fl] (e3)--(e4);
\draw[fl] (e4)--(e5); \draw[fl] (e5)--(e6); \draw[fl] (e6)--(e7);
\draw[fl] (e7)--(e8);
\node[font=\footnotesize\itshape, color=CouleurSecondaire,
    text width=3.5cm, align=left] at (5.5,-1.1)
    {Art.~52--59 Loi~2010\\perquisitions, saisies,\\réquisitions};
\node[font=\footnotesize\itshape, color=CouleurSecondaire,
    text width=3.5cm, align=left] at (5.5,-4.4)
    {Décret~69/DF/544\\expert agréé MJ};
\end{tikzpicture}
\caption{Procédure type d'investigation numérique au Cameroun}
\label{fig:procedure-cam}
\end{figure}

% ============================================================================
\section{Jurisprudence Camerounaise~: Premières Décisions de Référence}
\label{sec:jurisprudence}
% ============================================================================

La jurisprudence numérique camerounaise est naissante~: peu de décisions
publiées à ce jour. Trois affaires de référence~:

\begin{boxscenario}{Affaire CAMTEL c.~X (2018) --- première condamnation pour intrusion}
\textbf{Faits~:} intrusion non autorisée dans les systèmes de CAMTEL~;
modification de données de facturation au profit du prévenu.\\
\textbf{Qualification retenue~:} art.~65~al.~2 + art.~65~al.~3
Loi~2010/012 (accès illicite + aggravation pour modification de données).\\
\textbf{Preuves admises~:} logs système~; IP tracking~; analyse forensique
d'un expert agréé.\\
\textbf{Décision~:} 2~ans d'emprisonnement ferme (sursis refusé, art.~89)~;
5~millions FCFA d'amende~; dommages-intérêts à CAMTEL.\\
\textbf{Enseignement~:} première validation judiciaire des logs système
comme preuves numériques suffisantes. La chaîne de custody de l'expert
agréé était irréprochable.
\end{boxscenario}

\begin{boxscenario}{Affaire Ministère public c.~Y (2020) --- fraude Mobile Money}
\textbf{Faits~:} escroquerie via Mobile Money~; le prévenu interceptait
des codes OTP par SIM~swapping~; 47~victimes~; préjudice total~:
23~millions FCFA.\\
\textbf{Qualification retenue~:} art.~72 Loi~2010 (fraude informatique)
+ art.~318 Code pénal (escroquerie).\\
\textbf{Innovation~:} première utilisation des données USSD d'Orange
Cameroun comme preuve admissible.\\
\textbf{Enseignement~:} les données opérateurs (MTN MoMo, Orange Money)
constituent des preuves admissibles. Le cumul art.~72 Loi~2010 + art.~318
CP est validé par la jurisprudence.
\end{boxscenario}

\subsection{Défis de la jurisprudence camerounaise}

\begin{enumerate}
    \item \textbf{Formation des magistrats~:} insuffisante en technique~;
          le projet de formation continue des magistrats sur le droit du
          numérique est en cours au Ministère de la Justice.
    \item \textbf{Délais d'expertise~:} 1 à 3~mois en moyenne~;
          absence de laboratoire forensique national accrédité.
    \item \textbf{Jurisprudence non publiée~:} peu de décisions commentées
          et diffusées~; lacune que ce cours contribue à combler.
    \item \textbf{Absence de normes nationales d'accréditation~:}
          l'agrément des experts (Décret~69/DF/544) ne fixe pas de
          standard technique précis.
\end{enumerate}

\separateur

\begin{boxresume}
\textbf{Ce qu'il faut retenir de ce chapitre~:}
\begin{itemize}
    \item \textbf{Trois textes, toujours cités en cumul~:}
          \loicam{2010/012} (lex specialis cyber) + Code pénal
          (escroquerie, faux, diffamation) + \loicam{2024/017}
          (données personnelles).

    \item \textbf{Art.~89 Loi~2010~: sursis impossible.} Mentionner
          dans chaque rapport transmis au parquet.

    \item \textbf{L'adresse IP est une donnée personnelle} (art.~5
          Loi~2024). Toute intrusion consignant des IP active la Loi~2024
          en plus de la Loi~2010.

    \item \textbf{Art.~88 pour les téléphones chiffrés~:} refus de
          déverrouiller après réquisition judiciaire~= infraction
          autonome (1--5~ans). Ne jamais bloquer sur ce point.

    \item \textbf{Ransomware~:} art.~66--67 Loi~2010 + art.~63~al.~2
          Loi~2024 (peines doublées~: chiffrement~= atteinte
          à la disponibilité). + art.~22 Loi~2024 si la victime
          n'a pas notifié l'Autorité.

    \item \textbf{Personnes morales~:} Loi~2024 prévoit jusqu'à
          1~milliard FCFA (art.~71). Pour les infractions d'entreprises,
          toujours rechercher si la Loi~2024 est applicable~---
          l'effet dissuasif est sans commune mesure avec la Loi~2010.

    \item \textbf{Transfert illégal de données à l'étranger~:}
          art.~69 Loi~2024~: 3--10~ans. La sanction la plus lourde
          du dispositif national. L'autorisation préalable de l'Autorité
          de protection est obligatoire pour tout transfert.
\end{itemize}
\end{boxresume}

\bigskip

\begin{boxexo}[title={\ding{45}\quad Exercice~12.1 --- Qualification nationale \niveaubase}]
Pour chacun des 5~cas suivants, identifiez~: (A)~l'article Loi~2010/012
applicable, (B)~l'article Loi~2024/017 si pertinent, (C)~l'article Code
pénal en cumul~:
\begin{enumerate}[(a)]
    \item Un individu accède au compte Facebook d'une étudiante en
          utilisant son mot de passe dérobé par phishing et publie des
          photos intimes.
    \item Un employé copie la base de données clients de son employeur
          (50~000~entrées avec téléphone et email) et la vend à un concurrent.
    \item Un ransomware paralyse le système de facturation d'une clinique
          privée de Douala pendant 5~jours~; toutes les données patients
          sont chiffrées.
    \item Le directeur informatique d'une banque refuse de remettre la
          clé de déchiffrement du serveur de sauvegarde après réquisition
          formelle du juge d'instruction.
    \item Une startup camerounaise envoie sans autorisation les données
          personnelles de 200~000~utilisateurs vers ses serveurs AWS
          en Irlande.
\end{enumerate}
\end{boxexo}

\begin{boxexo}[title={\ding{45}\quad Exercice~12.2 --- Tableau de sanctions \niveaumoyen}]
En utilisant le Tableau~12.4~:
\begin{enumerate}[(a)]
    \item Quelle est la qualification qui expose l'auteur aux peines les
          plus sévères dans chaque cas~: le dirigeant d'une PME qui a
          vendu des données clients à l'étranger sans autorisation (1~000
          personnes concernées)~?
    \item Pour une entreprise qui a subi une fuite de données massive
          et n'a pas notifié l'Autorité de protection~: quelles infractions
          peut-on retenir contre l'entreprise elle-même (personne morale)~?
          Quelle est la sanction maximale potentielle~?
    \item Le sursis peut-il être accordé pour un condamné pour fraude
          informatique (art.~72 Loi~2010) ayant subtilisé 3~millions FCFA
          via Mobile Money~? Justifiez.
\end{enumerate}
\end{boxexo}

\begin{boxexo}[title={\ding{45}\quad Exercice~12.3 --- Simulation MOABI \niveauexpert}]
La Brigade de lutte contre la cybercriminalité de Yaoundé reçoit une plainte.
L'entreprise MOABI~SA (200~employés) signale~:
\begin{enumerate}
    \item Son site e-commerce a été piraté~: 45~000 numéros de cartes
          bancaires clients dérobés.
    \item Son serveur de paie a été chiffré par un ransomware
          (données RH inaccessibles depuis 4~jours).
    \item Elle n'a pas encore notifié ses clients ni l'Autorité de
          protection des données.
    \item Le suspect identifié a supprimé son téléphone principal
          mais possède un second téléphone chiffré qu'il refuse de
          déverrouiller.
\end{enumerate}
\begin{enumerate}[(a)]
    \item Listez toutes les infractions constituées avec les articles
          précis (distinguez~: infractions du suspect vs infractions de
          MOABI~SA).
    \item Identifiez les responsabilités de MOABI~SA en tant que personne
          morale~: quels articles de la Loi~2024 s'appliquent~?
    \item Rédigez un rapport de qualification transmissible au parquet
          (une page maximum)~: structure~: faits~/ qualifications~/ articles~/ sanctions encourues~/ mention art.~89.
    \item Quel est l'enjeu de l'art.~88 dans ce dossier~? Quelle démarche
          l'OPJ doit-il engager~?
\end{enumerate}
\end{boxexo}

\bigskip

\begin{boxinfo}
\textbf{Transition.} Le Chapitre~\ref{chap:deon} (Déontologie de
l'investigateur) clôt la Partie~III. Il aborde les obligations éthiques,
les conflits d'intérêt, les responsabilités personnelles et la conduite
de l'investigateur en situation difficile~--- la dimension humaine du
métier que les textes de loi ne couvrent pas.
\end{boxinfo}
